\theory{Information theory}{How can uncertainty, storage, and communication be measured, and how can messages survive compression and noise?}{Information theory treats information as the resolution of uncertainty. It gives limits for representing data, sending it through a noisy channel, and recovering it reliably.}{Data compression, error correction, telecommunication, statistics, cryptography, machine learning, neuroscience, and physics.}{Entropy, mutual information, and channel capacity quantify uncertainty, compression limits, and the effect of noise.}

\paragraph{History and the communication problem}
\bookfigure{information_coding_channel}{Information and coding theory follow a message through encoding, noise, and decoding.}
Claude Shannon formalized information theory in 1948 at Bell Labs. Harry Nyquist had related transmission speed to the number of signal levels, and Ralph Hartley had measured the number of distinguishable symbol sequences. Shannon's problem was simple to state: reproduce at one place a message selected at another, even when the channel adds noise \cite{information_theory_ref1}.

The theory separates three jobs. Source coding removes predictable repetition so a message takes fewer symbols. Channel coding adds carefully designed redundancy so the receiver can correct errors. Cryptographic coding hides information from an unwanted receiver. The mathematics is based on probability, statistics, and logarithms.

\end{historylist}
\section*{3. Theory}
\subsection*{Core theory}
\paragraph{Information and entropy}
If an event has probability $p$, its self-information is
\[
I(p)=-\log_2p.
\]
A rare event tells us more when it occurs. For a random variable $X$ with probabilities $p_i$, Shannon entropy is
\[
H(X)=-\sum_i p_i\log_2p_i.
\]
It measures average uncertainty before the outcome is known. A fair coin has entropy $1$ bit; a coin that always lands heads has entropy $0$. The binary entropy function
\[
H_b(p)=-p\log_2p-(1-p)\log_2(1-p)
\]
is largest at $p=1/2$.

The logarithm base determines units: bits or shannons for base two, nats for the natural logarithm, and hartleys for base ten. Entropy is largest when all possible symbols are equally likely. For independent variables, joint entropy adds; for dependent variables, knowing one reduces uncertainty about the other.

\paragraph{Conditional and mutual information}
Conditional entropy $H(X\mid Y)$ measures uncertainty about $X$ after $Y$ is known. Mutual information is
\[
I(X;Y)=H(X)-H(X\mid Y).
\]
It measures how much observing $Y$ teaches us about $X$. If two variables are independent, mutual information is zero. In a medical test, $X$ may be whether a disease is present and $Y$ a test result; mutual information measures the average reduction in diagnostic uncertainty.

The Kullback--Leibler divergence compares a proposed distribution with the true distribution:
\[
D(P\Vert Q)=\sum_xP(x)\log\frac{P(x)}{Q(x)}.
\]
It is nonnegative but not symmetric, so it is a divergence rather than an ordinary distance. Directed information measures information flowing from a past sequence to a future sequence and is useful for channels with memory.

\paragraph{Source coding and compression}
The source coding theorem says that a long sequence from a stationary source can be compressed to approximately its entropy rate bits per symbol, but not reliably below that rate. Huffman and arithmetic coding approach this limit for practical sources. ZIP files exploit repeated patterns; images, sound, and video use models of visual or auditory redundancy.

If a source has $n$ equally likely symbols, its entropy is $\log_2n$. If symbols are correlated, a compressor predicts the next symbol from the previous ones and stores only the surprise. Lossless compression reconstructs the exact data. Lossy compression permits a controlled distortion to achieve a much smaller representation.

\paragraph{Noisy channels and capacity}
A channel maps an input symbol to a random output. Its capacity is the greatest rate at which information can be transmitted with arbitrarily small error probability using long codes. Shannon's noisy-channel coding theorem says that rates below capacity are achievable, while rates above capacity are not \cite{information_theory_ref2}.

For a binary symmetric channel that flips a bit with probability $p$, the capacity is $1-H_b(p)$ bits per use. For a band-limited Gaussian channel with bandwidth $B$, signal power $S$, and noise power $N$, the Shannon--Hartley law gives
\[
C=B\log_2(1+S/N).
\]
Capacity is a limit, not a promise that every short practical code reaches it.

\paragraph{Codes, redundancy, and cryptography}
Error-correcting codes deliberately add redundancy. A repetition code sends 0 as 000 and 1 as 111; a majority vote corrects one error but uses bandwidth inefficiently. Hamming, Reed--Solomon, convolutional, turbo, and low-density parity-check codes approach capacity much more effectively. They made compact discs, mobile phones, DSL, satellite communication, and Voyager data possible \cite{information_theory_ref3}.

Cryptography uses codes and ciphers to control who can recover a message. Information-theoretic secrecy asks for security even against unlimited computation; the one-time pad achieves it when used correctly. Computational cryptography instead relies on hard problems such as discrete logarithms. Information theory also studies cryptanalysis, secrecy capacity, and leakage from side channels.

\paragraph{Applications across science}
In statistics, entropy and divergence measure model fit and guide inference. In machine learning, cross-entropy trains probabilistic predictors, while mutual information measures feature relevance. In neuroscience, information measures how a neural response reveals a stimulus. In bioinformatics it studies molecular codes and sequence evolution. In physics, thermodynamic entropy, black-hole entropy, Landauer's principle, and quantum information connect information to energy and matter.

Other applications include seismic exploration, signal processing, pattern recognition, anomaly detection, music and art analysis, imaging design, information retrieval, plagiarism detection, intelligence gathering, semiotics, and the study of biological and social networks \cite{information_theory_ref4}.

\paragraph{What the theory depends on and where it is useful}
Information theory depends on probability, statistics, logarithms, coding, and optimization. It helps coding theory design explicit codes, statistics quantify uncertainty, cryptography analyze secrecy, and machine learning compare distributions. Its model of information is mathematical and probabilistic; it does not by itself measure meaning, truth, or human understanding.

\section*{4. Important theorems and results}
\begin{enumerate}[leftmargin=*,itemsep=0.35em]
\item \textbf{Source coding theorem.} Entropy is the asymptotic lower limit for reliable lossless compression.

\item \textbf{Noisy-channel coding theorem.} Every rate below channel capacity can be transmitted with error tending to zero using suitable long codes; rates above capacity cannot.

\item \textbf{Data-processing inequality.} Processing data cannot increase the mutual information it contains about an original variable.

\item \textbf{Shannon--Hartley law.} Gaussian-channel capacity is $B\log_2(1+S/N)$.

\item \textbf{Fano's inequality.} The probability of decoding error bounds how much uncertainty remains, linking estimation accuracy to information \cite{information_theory_ref1}.
\end{enumerate}

\theoryapplications
\theoryconnections{information theory}{%
\item Probability supplies random variables and conditional distributions, while measure theory handles continuous sources.
\item Combinatorics and coding theory construct messages, and optimization studies the trade-off between rate, distortion, and reliability.
}{%
\item Information theory helps communication, compression, statistics, machine learning, cryptography, and physics quantify uncertainty and dependence.
\item Entropy and mutual information are useful only after the source, observation, and noise model have been specified.
}
\section*{Glossary}
\begin{description}[leftmargin=*,style=nextline,itemsep=0.3em]
\item[\textbf{Entropy.}] The average uncertainty of a random variable, measured in bits; it is maximized when all outcomes are equally likely.
\item[\textbf{Mutual information.}] The average reduction in uncertainty about one random variable obtained by observing another; it quantifies statistical dependence.
\item[\textbf{Channel capacity.}] The maximum rate at which information can be transmitted through a noisy channel with arbitrarily low error probability using suitable long codes.
\item[\textbf{Source coding.}] Compression of a data source to use as few bits as possible; the source coding theorem shows that the entropy rate is the fundamental limit for lossless compression.
\item[\textbf{Channel coding.}] The deliberate addition of structured redundancy to a message so that a receiver can detect and correct errors introduced by a noisy channel.
\item[\textbf{Kullback--Leibler divergence.}] A non-symmetric measure of how one probability distribution differs from another; it is nonnegative and zero only when the distributions agree.
\item[\textbf{Conditional entropy.}] $H(X\mid Y)$, the remaining uncertainty about $X$ after $Y$ is observed.
\item[\textbf{Self-information.}] $I(p)=-\log_2 p$, the information content of an event with probability $p$.
\item[\textbf{Binary symmetric channel.}] A channel that independently flips each input bit with probability $p$.
\item[\textbf{Shannon--Hartley law.}] The formula $C=B\log_2(1+S/N)$ for the capacity of a band-limited Gaussian channel.
\item[\textbf{Data-processing inequality.}] Processing data cannot increase mutual information about an original variable.
\end{description}
\section*{7. Sample Questions}
\emph{Exercise reference:} \cite{information_theory_ref1}. The cited edition is the source for the exercise set; consult its Exercises section for the official statement and numbering.
\begin{enumerate}[leftmargin=*,itemsep=0.9em]
\item \textbf{Exercise 1.} A random variable $X$ takes values in $\{a,b,c,d\}$ with probabilities $\frac{1}{2},\frac{1}{4},\frac{1}{8},\frac{1}{8}$. Compute $H(X)$.

\emph{Solution.} $H(X)=-\frac{1}{2}\log_2\frac{1}{2}-\frac{1}{4}\log_2\frac{1}{4}-\frac{1}{8}\log_2\frac{1}{8}-\frac{1}{8}\log_2\frac{1}{8}=\frac{1}{2}\cdot 1+\frac{1}{4}\cdot 2+\frac{1}{8}\cdot 3+\frac{1}{8}\cdot 3=\frac{1}{2}+\frac{1}{2}+\frac{3}{8}+\frac{3}{8}=\frac{7}{4}=1.75$ bits.

\item \textbf{Exercise 2.} Compute $H(X,Y)$ for two independent fair coins $X$ and $Y$.

\emph{Solution.} Since $X$ and $Y$ are independent with $H(X)=H(Y)=1$ bit, $H(X,Y)=H(X)+H(Y)=2$ bits. There are four equally likely outcomes, each with probability $1/4$, so $H(X,Y)=-4\cdot\frac{1}{4}\log_2\frac{1}{4}=2$ bits.

\item \textbf{Exercise 3.} A binary symmetric channel has crossover probability $p=0.25$. Compute its capacity in bits per use.

\emph{Solution.} $C=1-H_b(p)=1-H_b(0.25)=-0.25\log_2(0.25)-0.75\log_2(0.75)=0.5+0.75\cdot 0.41503\ldots=0.5+0.31127\ldots=0.8113$ bits. So $C\approx 0.1887$ bits per use.

\item \textbf{Exercise 4.} A source emits symbols from $\{A,B,C,D,E\}$ with equal probability. How many bits per symbol are needed for lossless compression, and how many bits would a naive fixed-length code use?

\emph{Solution.} Entropy: $H(X)=\log_2 5\approx 2.3219$ bits per symbol. A fixed-length code uses $\lceil\log_2 5\rceil=3$ bits per symbol. The source coding theorem says the optimal lossless code uses approximately $2.322$ bits per symbol, a savings of about $23\%$ over the fixed-length code.

\item \textbf{Exercise 5.} Two coins $X$ and $Y$ are not independent: $X$ is fair, and $Y=X$ with probability $0.9$ and $Y\neq X$ with probability $0.1$ (each with $0.05$ for the other value). Compute $I(X;Y)$.

\emph{Solution.} $H(X)=1$ bit. We compute $H(Y)$: $P(Y=H)=0.9\cdot 0.5+0.05\cdot 0.5=0.475$, $P(Y=T)=0.475$. So $H(Y)=-2(0.475)\log_2(0.475)\approx 0.9988$ bits. $H(X,Y)=\sum_{x,y}P(x,y)\log_2(1/P(x,y))$: $P(H,H)=0.45$, $P(T,T)=0.45$, $P(H,T)=0.05$, $P(T,H)=0.05$. So $H(X,Y)=-2(0.45)\log_2(0.45)-2(0.05)\log_2(0.05)\approx 0.9170+0.4322=1.3493$ bits. $I(X;Y)=H(X)+H(Y)-H(X,Y)=1+0.9988-1.3493\approx 0.6495$ bits.

\end{enumerate}
\section*{8. References}
\addcontentsline{toc}{section}{References}

\begin{thebibliography}{9}
\bibitem{information_theory_ref1} Claude E. Shannon, "A mathematical theory of communication," \emph{Bell System Technical Journal}, 27 (1948), 379--423 and 623--656.
\bibitem{information_theory_ref2} Thomas M. Cover and Joy A. Thomas, \emph{Elements of Information Theory}, 2nd ed., Wiley, 2006.
\bibitem{information_theory_ref3} David J. C. MacKay, \emph{Information Theory, Inference, and Learning Algorithms}, Cambridge University Press, 2003.
\bibitem{information_theory_ref4} Robert B. Ash, \emph{Information Theory}, Dover, 1990.
\end{thebibliography}
